\section{Основная матрица сравнений}

Центральная постановка работы строится вокруг четырех чисто монетарных вариантов денежно-кредитного правила: правила при полной информации; классического правила по оценённому состоянию; отобранного правила по оценённому состоянию; и отобранного правила по наблюдаемым переменным без явного этапа фильтрации. Такая матрица позволяет отдельно оценить ценность информации, выигрыш от оценивания скрытого состояния и выигрыш от более гибкой формы правила.

По четырем базовым сценариям --- базовый макроэкономический набор и ограниченный информационный набор, каждый в сочетании с умеренной или высокой разделимостью режимов --- отобранное правило по оценённому состоянию уменьшает накопленную функцию потерь относительно классического правила по оценённому состоянию на 56.7--71.1\%. При этом выигрыш от оценивания скрытого состояния зависит от информационного режима. В сценариях с базовым макроэкономическим набором правило по наблюдаемым переменным и правило по оценённому состоянию почти не различаются, а при высокой разделимости режимов прямое правило по наблюдаемым данным даже слегка лучше. Напротив, при ограниченном информационном наборе преимущество правила по оценённому состоянию становится выраженным: относительный выигрыш по сравнению с правилом, использующим только наблюдаемые переменные, достигает 50.9\%, а в сценарии ``ограниченный информационный набор × высокая разделимость режимов'' доверительный интервал для разницы потерь целиком лежит ниже нуля в пользу фильтрации.

Правило при полной информации используется здесь не как глобальная верхняя граница, а как внутренний ориентир внутри одной и той же семьи правил. Сопоставление с ним показывает, насколько дорого обходится потеря информации и насколько этот разрыв удается сократить за счет более качественного представления состояния. Тем самым главный результат этапа 6 связан не с усложнением пространства действий, а с тем, какую информацию регулятор действительно получает и как она переводится в решение по ставке.

\begin{table}[htbp]
\centering
\small
\caption{Основная матрица сравнений денежно-кредитных правил}
\label{tab:stage6_core_main_matrix}
\begin{tabular}{lcccc}
\toprule
Сценарий & Полная информация & Классическое по оценённому состоянию & Отобранное по оценённому состоянию & Отобранное по наблюдаемым переменным \\
\midrule
Базовый макроэкономический набор × умеренная разделимость режимов & 0.0061 & 0.0076 & 0.0022 & 0.0020 \\
Базовый макроэкономический набор × высокая разделимость режимов & 0.0070 & 0.0085 & 0.0026 & 0.0024 \\
Ограниченный информационный набор × умеренная разделимость режимов & 0.0061 & 0.0085 & 0.0037 & 0.0042 \\
Ограниченный информационный набор × высокая разделимость режимов & 0.0070 & 0.0102 & 0.0036 & 0.0077 \\
\bottomrule
\end{tabular}
\end{table}

\begin{table}[htbp]
\centering
\small
\caption{Выигрыш от оценивания скрытого состояния, гибкости правила и разрыв до полной информации}
\label{tab:stage6_core_values}
\begin{tabular}{lccc}
\toprule
Сценарий & Выигрыш от оценивания скрытого состояния, \% & Выигрыш от гибкости правила, \% & Разрыв до полной информации \\
\midrule
Базовый макроэкономический набор × умеренная разделимость режимов & -9.4410 & 71.0806 & -0.0039 \\
Базовый макроэкономический набор × высокая разделимость режимов & -8.5734 & 69.3113 & -0.0044 \\
Ограниченный информационный набор × умеренная разделимость режимов & 11.9298 & 56.6728 & -0.0024 \\
Ограниченный информационный набор × высокая разделимость режимов & 52.5890 & 64.2882 & -0.0034 \\
\bottomrule
\end{tabular}
\end{table}
